\subsection{Priors, Posteriors, and Conjugacy}

Bayesian inference treats unknown parameters as uncertain quantities and updates a prior distribution by the likelihood supplied by the data.

\begin{definition} \label{def:prior-posterior-predictive} \label{def:conjugate-prior-family} \hlt{Prior, Posterior, Predictive, and Conjugacy}\\
Let $\Pi$ be a prior on $\Theta$ and $p_\theta(x)$ a dominated likelihood.
\begin{enumerate}[label=(\roman*)]
\setlength{\itemsep}{0pt}
\item \hlt{Prior}: $\Pi$ encodes uncertainty about $\theta\in\Theta$ before observing $x$.
\item \hlt{Evidence}: $m(x)\defeq\int_\Theta p_\vartheta(x)\,d\Pi(\vartheta)$, assumed finite and positive.
\item \hlt{Posterior}: $d\Pi(\theta\mid x)\defeq \frac{p_\theta(x)d\Pi(\theta)}{m(x)}$.
\item \hlt{Predictive}: for future $X_{\mathrm{new}}$, $\mathbb{P}(X_{\mathrm{new}}\in A\mid x)=\int_\Theta P_\theta(A)\,d\Pi(\theta\mid x)$.
\item \hlt{Conjugacy}: $\{\Pi_\eta:\eta\in H\}$ is conjugate if $\Pi_\eta(\cdot\mid x)=\Pi_{\eta(x)}$ for some $\eta(x)\in H$.
\end{enumerate}
\end{definition}

\begin{example} \label{ex:beta-binomial-conjugacy} \hlt{Beta--Binomial Updating}\\
If $X\mid p\sim\operatorname{Binomial}(n,p)$ and $p\sim\operatorname{Beta}(a,b)$, then $p\mid X=x\sim\operatorname{Beta}(a+x,b+n-x)$.\\
Thus $\mathbb{E}[p\mid x]=\frac{a+x}{a+b+n}$, a weighted average of $\frac{a}{a+b}$ and $\frac{x}{n}$.
\end{example}

\begin{example} \label{ex:gamma-poisson-conjugacy-rate} \hlt{Gamma--Poisson Updating}\\
Let $X_1,\ldots,X_n\mid\lambda$ be iid $\operatorname{Poisson}(\lambda)$ and $\lambda\sim\operatorname{Gamma}(\alpha,\nu)$ in rate form.\\
Writing $s\defeq\sum\limits_{i=1}^nx_i$, the posterior kernel is $\pi(\lambda\mid x)\propto\lambda^{\alpha+s-1}e^{-(\nu+n)\lambda}$, hence $\lambda\mid x\sim\operatorname{Gamma}(\alpha+s,\nu+n)$.\\
Thus $\mathbb{E}[\lambda\mid x]=\frac{\alpha+s}{\nu+n}$, $\operatorname{Var}(\lambda\mid x)=\frac{\alpha+s}{(\nu+n)^2}$, $\operatorname{mode}(\lambda\mid x)=\frac{\alpha+s-1}{\nu+n}$ if $\alpha+s>1$.
\end{example}

\begin{example} \label{ex:normal-normal-conjugacy} \hlt{Normal Mean with Normal Prior}\\
Let $X_1,\ldots,X_n\mid\mu$ be iid $N(\mu,\sigma^2)$ with known $\sigma^2$, and let $\mu\sim N(m_0,\tau_0^2)$.\\
Then $\mu\mid X\sim N(m_n,\tau_n^2)$, where $\tau_n^{-2}=\tau_0^{-2}+n\sigma^{-2}$ and $m_n=\tau_n^2(\tau_0^{-2}m_0+n\sigma^{-2}\overline{X})$.
\end{example}

\begin{example} \label{ex:normal-precision-gamma-conjugacy} \hlt{Known Mean and Unknown Precision}\\
Let $X_1,\ldots,X_n\mid\xi$ be iid $N(\theta_0,\xi^{-1})$ with known $\theta_0$, and let $\xi\sim\operatorname{Gamma}(\alpha,\lambda)$ in rate form.\\
Writing $q\defeq\sum\limits_{i=1}^n(x_i-\theta_0)^2$, the posterior is $\xi\mid x\sim\operatorname{Gamma}(\alpha+\frac{n}{2},\lambda+\frac{q}{2})$.\\
With a flat gamma prior, $\sigma^2=\xi^{-1}$ centers near the MLE $\frac{q}{n}$.
\end{example}

\begin{definition} \label{rem:improper-priors-need-proper-posteriors} \hlt{Improper Prior}\\
An \hlt{improper prior} is a $\sigma$-finite prior measure $\Pi$ on $\Theta$ with $\Pi(\Theta)=\infty$.
For likelihood $p_\theta(x)$, it is usable only when $0<m_\Pi(x)\defeq\int_\Theta p_\theta(x)\,d\Pi(\theta)<\infty$, in which case $d\Pi(\theta\mid x)\defeq\frac{p_\theta(x)d\Pi(\theta)}{m_\Pi(x)}$ is a proper posterior.
\end{definition}

\begin{remark} \label{rem:bayesian-conditioning-measure-theory} \hlt{Conditioning Viewpoint}\\
Bayesian updating is conditioning under the joint law $dP(\theta,x)=p_\theta(x)\,d\Pi(\theta)\,dx$; the posterior $\Pi(\cdot\mid x)$ is a regular conditional law of $\theta$ given $X=x$.
\end{remark}
